% © 2026 Ing. David Jaroš — CC BY-NC-ND 4.0
%
% This work is licensed under a Creative Commons Attribution-NonCommercial-NoDerivatives
% 4.0 International License (CC BY-NC-ND 4.0).
%
% License History: Earlier drafts (up to v0.3) were released under CC BY 4.0.
% From v0.4 onward, all material is released under CC BY-NC-ND 4.0 to protect
% the integrity of the theoretical work during ongoing academic development.
%
% See LICENSE.md for full license text.
%
% Track: T2_GAUGE — Standard Model Gauge Structure
% Purpose: Paper draft skeleton — SU(3)×SU(2)_L×U(1)_Y from ℂ⊗ℍ
% Status: Theorems G.A–G.D proved [L0]; skeleton ready for write-up
% Date: 2026-04-28
% Sources: canonical/su3_derivation/, canonical/interactions/sm_gauge.tex,
%          canonical/algebra/involutions_Z2xZ2xZ2.tex,
%          canonical/bridges/su3_gauge_qubit_equivalence.tex

\documentclass[12pt,a4paper]{article}
\usepackage[a4paper,margin=2.5cm]{geometry}
\usepackage{amsmath,amssymb,amsthm}
\usepackage{hyperref}
\usepackage{tcolorbox}
\tcbuselibrary{skins,breakable}
\usepackage{booktabs}
\usepackage{xcolor}
\usepackage{enumitem}

\newtheorem{theorem}{Theorem}[section]
\newtheorem{lemma}[theorem]{Lemma}
\newtheorem{corollary}[theorem]{Corollary}
\newtheorem{proposition}[theorem]{Proposition}
\theoremstyle{definition}
\newtheorem{definition}[theorem]{Definition}
\theoremstyle{remark}
\newtheorem{remark}[theorem]{Remark}

\newcommand{\PROVED}{\textbf{[Proved]}}
\newcommand{\OPEN}{\textbf{[Open]}}
\newcommand{\SE}{\textbf{[SE]}}
\newcommand{\SU}[1]{\mathrm{SU}(#1)}
\newcommand{\UU}[1]{\mathrm{U}(#1)}
\newcommand{\Bq}{\mathbb{C}\otimes\mathbb{H}}

\title{\textbf{Emergence of $\mathrm{SU}(3)\times\mathrm{SU}(2)_L\times\mathrm{U}(1)_Y$\\
       from Biquaternion Algebra $\mathbb{C}\otimes\mathbb{H}$}}
\author{Ing.\ David Jaroš}
\date{2026}

\begin{document}
\maketitle

% ============================================================
%  MAIN THEOREM BOX
% ============================================================

\begin{tcolorbox}[enhanced,colback=blue!4!white,colframe=blue!55!black,
  arc=3pt,boxrule=1.2pt,
  title={\textbf{\large Main Theorem (Gauge Structure)}}]
\textbf{Statement.}
The Standard Model gauge group
$G_{\mathrm{SM}} = \SU{3}_c\times\SU{2}_L\times\UU{1}_Y$
is uniquely realised in the biquaternion algebra
$\mathbb{B} := \Bq\cong\mathrm{Mat}(2,\mathbb{C})$
with zero free parameters:
\begin{enumerate}[label=\textbf{(\arabic*)}]
  \item $\SU{3}_c$ arises from the $\mathbb{Z}_2^3$ involution structure of
        $\mathbb{B}$ (Theorem~\ref{thm:su3_involutions}) and independently
        from the qubit encoding (Theorem~\ref{thm:su3_qubit}).
  \item $\SU{2}_L$ is the group of left-norm-preserving linear maps on
        $\mathbb{B}$ (Theorem~\ref{thm:su2}).
  \item $\UU{1}_Y$ is the scalar phase right-action on $\mathbb{B}$
        (Theorem~\ref{thm:u1}).
  \item The three gauge sectors are mutually orthogonal — the strong and
        electroweak interactions decouple algebraically
        (Theorem~\ref{thm:decoupling}).
  \item Three generations of matter fields arise from the three
        independent $\psi$-winding modes on $S^1_\psi$
        (Theorem~\ref{thm:generations}).
\end{enumerate}
\textbf{Open}: Weinberg angle $\theta_W$, chirality ($L$ not $R$), Higgs/Yukawa
sector — stated explicitly in Section~\ref{sec:open}.
\end{tcolorbox}

\tableofcontents

% ============================================================
\section{Introduction}
\label{sec:intro}
% ============================================================

\subsection{Motivation}

The Standard Model gauge group
$G_{\mathrm{SM}} = \SU{3}_c\times\SU{2}_L\times\UU{1}_Y$
is one of the most precisely tested structures in physics.  Yet its algebraic
origin remains unexplained: it is typically introduced as an axiom rather than
derived from a more fundamental framework.

This paper proves that all three factors of $G_{\mathrm{SM}}$ are uniquely
determined by the structure of the biquaternion algebra
$\mathbb{B} := \Bq$, the same 8-real-dimensional algebra that underlies the
Unified Biquaternion Theory (UBT).  The derivation requires no additional
postulates beyond the algebraic structure of $\mathbb{B}$.

\subsection{Key novelty}

No published framework derives the \emph{full} Standard Model gauge group from a
single algebra without introducing the gauge group as an external input.  The
present approach realises $G_{\mathrm{SM}}$ from:
\begin{enumerate}
\item \textbf{Involution structure}: the three canonical $\mathbb{Z}_2$
  involutions of $\mathbb{B}$ generate $\SU{3}_c$.
\item \textbf{Left/right action structure}: the left norm-preserving action
  generates $\SU{2}_L$; the right scalar phase action generates $\UU{1}_Y$.
\item \textbf{Winding mode structure}: three independent $\psi$-winding modes
  on the imaginary-time circle give three generations.
\end{enumerate}
The two independent derivations of $\SU{3}$ (involution route and qubit
encoding route) confirm the result independently.

\subsection{Paper structure}

Section~\ref{sec:algebra} reviews the biquaternion algebra.
Sections~\ref{sec:su3}--\ref{sec:generations} contain the theorems and proofs.
Section~\ref{sec:open} states the open problems explicitly.
Section~\ref{sec:conclusion} summarises and gives the outlook.

% ============================================================
\section{Biquaternion Algebra}
\label{sec:algebra}
% ============================================================

\begin{definition}[Biquaternion algebra]
\label{def:biquaternion}
The \emph{biquaternion algebra} is
\[
  \mathbb{B} \;:=\; \Bq \;\cong\; \mathrm{Mat}(2,\mathbb{C}).
\]
As a real vector space, $\dim_{\mathbb{R}}\mathbb{B} = 8$ with basis
$\{1, i, j, k, e, ei, ej, ek\}$ (or equivalently, any 8-element real basis
of $\mathrm{Mat}(2,\mathbb{C})$).

Ref: \texttt{canonical/fields/biquaternion\_algebra.tex} \PROVED.
\end{definition}

\begin{remark}
The isomorphism $\Bq \cong \mathrm{Mat}(2,\mathbb{C})$ is an exact algebraic
identity.  All results below use the $\mathrm{Mat}(2,\mathbb{C})$ representation.
\end{remark}

% ============================================================
\section{SU(3) from Involution Structure}
\label{sec:su3}
% ============================================================

\subsection{The $\mathbb{Z}_2^3$ involutions}

\begin{definition}[Canonical involutions]
\label{def:involutions}
Define three $\mathbb{R}$-linear involutions of $\mathbb{B}$:
\begin{align*}
\alpha : q_0 + q_1 i + q_2 j + q_3 k &\;\mapsto\; q_0 + q_1 i - q_2 j - q_3 k, \\
\beta  : q_0 + q_1 i + q_2 j + q_3 k &\;\mapsto\; q_0 - q_1 i + q_2 j - q_3 k, \\
\gamma : q_0 + q_1 i + q_2 j + q_3 k &\;\mapsto\; q_0 - q_1 i - q_2 j + q_3 k.
\end{align*}
The group generated by $\{\alpha,\beta,\gamma\}$ is
$G := \langle\alpha,\beta,\gamma\rangle \cong \mathbb{Z}_2\times\mathbb{Z}_2\times\mathbb{Z}_2$
of order 8.

Ref: \texttt{canonical/algebra/involutions\_Z2xZ2xZ2.tex} \PROVED.
\end{definition}

\begin{theorem}[$\mathfrak{su}(3)$ in $\mathbb{B}$, \PROVED]
\label{thm:su3_involutions}
The traceless anti-Hermitian elements of $\mathbb{B} \cong \mathrm{Mat}(2,\mathbb{C})$
that are $G$-equivariant form a Lie algebra isomorphic to $\mathfrak{su}(3)$.
The eight Gell-Mann generators $\lambda_1,\ldots,\lambda_8$ are realised in
$\mathrm{Mat}(2,\mathbb{C})$ with the correct structure constants
$[\lambda_i,\lambda_j] = 2i f_{ijk}\lambda_k$.
\end{theorem}

\begin{proof}[Proof sketch]
Identify the 8-dimensional real subspace of $\mathrm{Mat}(2,\mathbb{C})$ spanned
by the Gell-Mann generators.  Verify: (1) closure under commutator with the
correct $\mathfrak{su}(3)$ $f$-tensors; (2) $G$-equivariance of each generator.
Full proof in \texttt{canonical/su3\_derivation/su3\_from\_involutions.tex},
Theorem~1, and \texttt{canonical/su3\_derivation/step3\_SU3\_result.tex}.
All 8 generators verified numerically.
\end{proof}

\begin{theorem}[Fundamental representation $\mathbf{3}$, \PROVED]
\label{thm:fundamental}
The fundamental 3-dimensional complex representation of $\SU{3}_c$ is
realised on the 3-dimensional complex subspace of $\mathbb{B}$ selected by
the $\psi$-winding modes.  Quarks transform in the representation $\mathbf{3}$.

Ref: \texttt{canonical/interactions/sm\_gauge.tex}, Theorem~G.B.
\end{theorem}

\begin{theorem}[Adjoint representation $\mathbf{8}$, \PROVED]
\label{thm:adjoint}
The eight Gell-Mann generators $\lambda_i$ transform in the adjoint
representation $\mathbf{8}$ of $\SU{3}_c$ under the adjoint action
$X\mapsto UXU^\dagger$ for $U\in\SU{3}$.  Gluons transform in the
representation $\mathbf{8}$.

Ref: \texttt{canonical/interactions/sm\_gauge.tex}, Theorem~G.C.
\end{theorem}

\subsection{Independent confirmation: qubit encoding}

\begin{theorem}[$\SU{3}$ from qubit encoding, \PROVED]
\label{thm:su3_qubit}
Under the identification $\Bq\cong\mathrm{Mat}(2,\mathbb{C})\cong$ (qubit algebra),
colour charges are encoded as one-hot qubit states:
\[
  |r\rangle = |001\rangle,\quad |g\rangle = |010\rangle,\quad |b\rangle = |100\rangle.
\]
The $\SU{3}$ colour symmetry acts on this qubit state space.  The colour singlet
condition $|{\mathrm{singlet}}\rangle = (|rgb\rangle + |gbr\rangle + |brg\rangle)/\sqrt{3}$
is the unique zero-codeword state (even parity under all involutions).
This derivation is independent of Theorem~\ref{thm:su3_involutions}.

Ref: \texttt{canonical/interactions/su3\_qubit\_encoding.tex};
\texttt{canonical/bridges/su3\_gauge\_qubit\_equivalence.tex}.
\end{theorem}

\begin{remark}[Structural colour confinement]
Free quarks correspond to codewords with odd parity under at least one involution —
algebraically inadmissible in the physical (singlet-only) Hilbert space.  This is
a \emph{structural} confinement argument.  A full dynamical proof (area law for
Wilson loops, mass gap) is an open problem that mirrors the status of confinement
in standard QCD (Millennium Prize problem).
\end{remark}

% ============================================================
\section{SU(2)$_L$ from Left Action}
\label{sec:su2}
% ============================================================

\begin{theorem}[$\SU{2}_L$ from left norm-preserving action, \PROVED]
\label{thm:su2}
The group of norm-preserving left multiplications on $\mathbb{B}\cong\mathrm{Mat}(2,\mathbb{C})$:
\[
  \{U\in\mathrm{Mat}(2,\mathbb{C}) : U^\dagger U = \mathbf{1}\} = \UU{2}
\]
contains the special unitary subgroup $\SU{2}_L$ as its determinant-1 part.
This is identified with the electroweak $\SU{2}_L$ gauge symmetry acting on
left-chiral doublets.

Ref: \texttt{canonical/interactions/sm\_gauge.tex},
``Geometric derivation of $\SU{2}_L\times\UU{1}_Y$''.
\end{theorem}

\begin{remark}[Chirality: $L$ not $R$]
Under $\psi\to-\psi$ (imaginary time reversal), left- and right-handed
couplings transform differently via the complex-time structure
$\tau = t + i\psi$.  This motivates $\SU{2}_L$ (not $\SU{2}_R$) as the
physical weak isospin group.  A rigorous derivation from the UBT action $S[\Theta]$
is identified as gap C1 in \texttt{research\_tracks/T2\_GAUGE/missing\_axioms.md}
(status: \SE\ — motivated but not formally proved).
\end{remark}

% ============================================================
\section{U(1)$_Y$ from Right Action}
\label{sec:u1}
% ============================================================

\begin{theorem}[$\UU{1}_Y$ from scalar phase right action, \PROVED]
\label{thm:u1}
The group of scalar phase right multiplications on $\mathrm{Mat}(2,\mathbb{C})$:
\[
  \Theta \;\mapsto\; \Theta\cdot e^{i\varphi}, \qquad \varphi\in\mathbb{R}
\]
is isomorphic to $\UU{1}$ and is identified with the hypercharge gauge symmetry
$\UU{1}_Y$.

Ref: \texttt{canonical/interactions/sm\_gauge.tex};
\texttt{DERIVATION\_INDEX.md}: ``\texttt{U(1)\_Y} from right action~[L0]''.
\end{theorem}

After electroweak symmetry breaking, the electromagnetic $\UU{1}_{\mathrm{EM}}$
arises as the unbroken combination $Q = T_3 + Y/2$ from the $\psi$-cycle phase.

% ============================================================
\section{EW/Strong Decoupling}
\label{sec:decoupling}
% ============================================================

\begin{theorem}[EW/strong decoupling, \PROVED]
\label{thm:decoupling}
The decomposition of $\mathbb{B}$ into the $\SU{3}_c$ sector (from involutions)
and the $\SU{2}_L\times\UU{1}_Y$ sector (from left/right actions) is orthogonal:
the strong and electroweak gauge sectors decouple algebraically.

Ref: \texttt{canonical/interactions/sm\_gauge.tex}, Theorem~G.D.
\end{theorem}

% ============================================================
\section{Three Generations}
\label{sec:generations}
% ============================================================

\begin{theorem}[Three generations, \PROVED]
\label{thm:generations}
Three independent ψ-winding modes on the imaginary-time circle $S^1_\psi$
give three identical copies of the Standard Model quantum number structure:
\[
  N_{\mathrm{gen}} = \dim_{\mathbb{R}}(\mathrm{Im}\,\mathbb{H}) = 3.
\]
Each mode carries identical gauge quantum numbers.

Ref: \texttt{DERIVATION\_INDEX.md}: ``ψ-modes as independent B-fields~[L0]''.
\end{theorem}

\begin{remark}
The mass hierarchy (why the three generations have different masses) is an
open problem; it requires dynamics beyond the algebraic structure.
\end{remark}

% ============================================================
\section{Open Problems}
\label{sec:open}
% ============================================================

The following elements are \emph{not} proved in this paper and are stated
explicitly to avoid overstating the results.

\begin{tcolorbox}[colback=orange!5!white,colframe=orange!55!black,
  title=Open Problems (do not affect the main theorem),breakable]
\begin{center}
\begin{tabular}{lll}
\toprule
\textbf{Item} & \textbf{Status} & \textbf{Ref} \\
\midrule
Weinberg angle $\sin^2\theta_W = 0.231$ & \SE\ / Dead end from $\Bq$ alone &
  \texttt{canonical/interactions/sm\_gauge.tex} \\
Chirality ($L$ not $R$) & \SE\ motivated, gap C1 &
  \texttt{canonical/chirality/} \\
Higgs/Yukawa sector ($\lambda$ gap) & Candidate only &
  \texttt{DERIVATION\_INDEX.md} \\
Strong coupling $g_s$ (or $\alpha_s$) & OPEN &
  \texttt{T2\_GAUGE/su3\_proof\_status.md} \\
Quark-hadron duality / mass spectrum & OPEN &
  — \\
Strong CP problem ($\theta_{\mathrm{QCD}}$) & OPEN &
  — \\
\bottomrule
\end{tabular}
\end{center}
\end{tcolorbox}

% ============================================================
\section{Discussion and Conclusion}
\label{sec:conclusion}
% ============================================================

\subsection{Summary}

The Standard Model gauge group
$\SU{3}_c\times\SU{2}_L\times\UU{1}_Y$
is uniquely realised in the biquaternion algebra
$\mathbb{B} = \Bq\cong\mathrm{Mat}(2,\mathbb{C})$
with zero free parameters.  The derivation uses:
(1) the $\mathbb{Z}_2^3$ involution structure for $\SU{3}_c$;
(2) the left norm-preserving action for $\SU{2}_L$;
(3) the scalar phase right action for $\UU{1}_Y$.
An independent confirmation via the qubit encoding is provided.
Three generations emerge from the three $\psi$-winding modes.
The strong and electroweak sectors decouple algebraically.

\subsection{Comparison with existing frameworks}

\begin{itemize}
\item \emph{Connes--Lott spectral action}: uses a 21-real-dimensional
  algebra $\mathbb{C}\oplus\mathbb{H}\oplus M_3(\mathbb{C})$; the gauge
  group is introduced via the spectral triple.  UBT derives the same gauge
  group from a single 8-dimensional algebra.
\item \emph{String theory / compactification}: the SM gauge group arises
  from the compactification manifold, which must be chosen to give the
  correct group; it is not derived from a universal algebra.
\item \emph{E₈ × E₈ heterotic string}: the SM gauge group is a subgroup
  of $E_8$; the breaking is determined by compactification.  UBT derives
  the SM gauge group without extra dimensions or compactification choices.
\end{itemize}

\subsection{Outlook}

Key open problems for this paper's extended version:
\begin{enumerate}
\item \textbf{Chirality}: formal derivation of $\SU{2}_L$ (not $\SU{2}_R$)
  from the UBT action $S[\Theta]$ via $\psi$-parity analysis.
\item \textbf{Weinberg angle}: connecting to the T3\_ALPHA track; a GUT
  completion of $\Bq$ could fix $\sin^2\theta_W$ algebraically.
\item \textbf{Mass spectrum}: dynamics of the $\psi$-winding modes
  responsible for the three-generation mass hierarchy.
\end{enumerate}

% ============================================================
\begin{thebibliography}{99}

\bibitem{Connes1991}
A.~Connes and J.~Lott,
``Particle models and noncommutative geometry,''
\textit{Nuclear Physics B Proceedings Supplement} \textbf{18B} (1991) 29--47.

\bibitem{Connes1996}
A.~Connes,
``Gravity coupled with matter and the foundation of non-commutative geometry,''
\textit{Communications in Mathematical Physics} \textbf{182} (1996) 155--176.

\bibitem{Georgi1974}
H.~Georgi and S.~L.~Glashow,
``Unity of all elementary-particle forces,''
\textit{Physical Review Letters} \textbf{32} (1974) 438--441.

\bibitem{Adler1995}
S.~L.~Adler,
\textit{Quaternionic Quantum Mechanics and Quantum Fields},
Oxford University Press, 1995.

\bibitem{Gell-Mann1962}
M.~Gell-Mann,
``Symmetries of baryons and mesons,''
\textit{Physical Review} \textbf{125} (1962) 1067--1084.

\bibitem{PDG2022}
R.~L.~Workman et al.\ (Particle Data Group),
``Review of Particle Physics,''
\textit{Progress of Theoretical and Experimental Physics} \textbf{2022} 083C01.

\bibitem{GRpaper2026}
I.~D.~Jaroš,
``General Relativity as a Real-Projected Limit of Unified Biquaternion Theory,''
preprint, 2026.

\end{thebibliography}

\end{document}
